% ---------------------------------------------------------------------
% RECUPERO (file R1, R2) - conversione decimale -> binario
%   fila R1: 179=10110011, 427=110101011, 302=100101110, 871=1101100111
%   fila R2: 478=111011110, 203=11001011, 941=1110101101, 349=101011101
%   4 numeri a 3 cifre (8, 9, 9, 10 bit), un solo pari; numeri diversi dalle file A, B, C
% ---------------------------------------------------------------------
\begin{esercizio}[8]
Converti in binario:
\RR{%
\par\smallskip\noindent\hspace*{1em}\begin{tabular}{@{}p{0.42\linewidth}p{0.42\linewidth}@{}}
  \textbf{a)}~$179_{10} = \ldots\ldots\ldots_2$ & \textbf{b)}~$427_{10} = \ldots\ldots\ldots_2$\\[4pt]
  \textbf{c)}~$302_{10} = \ldots\ldots\ldots_2$ & \textbf{d)}~$871_{10} = \ldots\ldots\ldots_2$
\end{tabular}}{%
\par\smallskip\noindent\hspace*{1em}\begin{tabular}{@{}p{0.42\linewidth}p{0.42\linewidth}@{}}
  \textbf{a)}~$478_{10} = \ldots\ldots\ldots_2$ & \textbf{b)}~$203_{10} = \ldots\ldots\ldots_2$\\[4pt]
  \textbf{c)}~$941_{10} = \ldots\ldots\ldots_2$ & \textbf{d)}~$349_{10} = \ldots\ldots\ldots_2$
\end{tabular}}

\begin{solution}
\RR{%
\noindent
  \begin{minipage}[t]{0.24\linewidth}
  \textbf{a)}\\[2pt]
  {\small\renewcommand{\arraystretch}{1.0}\begin{tabular}{r|l}
    179 & 1\\
    89 & 1\\
    44 & 0\\
    22 & 0\\
    11 & 1\\
    5 & 1\\
    2 & 0\\
    1 & 1\\
    0 & \\
  \end{tabular}}
  \end{minipage}\hfill
  \begin{minipage}[t]{0.24\linewidth}
  \textbf{b)}\\[2pt]
  {\small\renewcommand{\arraystretch}{1.0}\begin{tabular}{r|l}
    427 & 1\\
    213 & 1\\
    106 & 0\\
    53 & 1\\
    26 & 0\\
    13 & 1\\
    6 & 0\\
    3 & 1\\
    1 & 1\\
    0 & \\
  \end{tabular}}
  \end{minipage}\hfill
  \begin{minipage}[t]{0.24\linewidth}
  \textbf{c)}\\[2pt]
  {\small\renewcommand{\arraystretch}{1.0}\begin{tabular}{r|l}
    302 & 0\\
    151 & 1\\
    75 & 1\\
    37 & 1\\
    18 & 0\\
    9 & 1\\
    4 & 0\\
    2 & 0\\
    1 & 1\\
    0 & \\
  \end{tabular}}
  \end{minipage}\hfill
  \begin{minipage}[t]{0.24\linewidth}
  \textbf{d)}\\[2pt]
  {\small\renewcommand{\arraystretch}{1.0}\begin{tabular}{r|l}
    871 & 1\\
    435 & 1\\
    217 & 1\\
    108 & 0\\
    54 & 0\\
    27 & 1\\
    13 & 1\\
    6 & 0\\
    3 & 1\\
    1 & 1\\
    0 & \\
  \end{tabular}}
  \end{minipage}\\[6pt]
\par\smallskip\noindent\hspace*{1em}\begin{tabular}{@{}p{0.42\linewidth}p{0.42\linewidth}@{}}
  \textbf{a)}~$179_{10} = 10110011_2$ & \textbf{b)}~$427_{10} = 110101011_2$\\[4pt]
  \textbf{c)}~$302_{10} = 100101110_2$ & \textbf{d)}~$871_{10} = 1101100111_2$
\end{tabular}\par\smallskip
  Si divide per 2 fino ad arrivare a 0 e si leggono i resti dal basso verso l'alto.}{%
\noindent
  \begin{minipage}[t]{0.24\linewidth}
  \textbf{a)}\\[2pt]
  {\small\renewcommand{\arraystretch}{1.0}\begin{tabular}{r|l}
    478 & 0\\
    239 & 1\\
    119 & 1\\
    59 & 1\\
    29 & 1\\
    14 & 0\\
    7 & 1\\
    3 & 1\\
    1 & 1\\
    0 & \\
  \end{tabular}}
  \end{minipage}\hfill
  \begin{minipage}[t]{0.24\linewidth}
  \textbf{b)}\\[2pt]
  {\small\renewcommand{\arraystretch}{1.0}\begin{tabular}{r|l}
    203 & 1\\
    101 & 1\\
    50 & 0\\
    25 & 1\\
    12 & 0\\
    6 & 0\\
    3 & 1\\
    1 & 1\\
    0 & \\
  \end{tabular}}
  \end{minipage}\hfill
  \begin{minipage}[t]{0.24\linewidth}
  \textbf{c)}\\[2pt]
  {\small\renewcommand{\arraystretch}{1.0}\begin{tabular}{r|l}
    941 & 1\\
    470 & 0\\
    235 & 1\\
    117 & 1\\
    58 & 0\\
    29 & 1\\
    14 & 0\\
    7 & 1\\
    3 & 1\\
    1 & 1\\
    0 & \\
  \end{tabular}}
  \end{minipage}\hfill
  \begin{minipage}[t]{0.24\linewidth}
  \textbf{d)}\\[2pt]
  {\small\renewcommand{\arraystretch}{1.0}\begin{tabular}{r|l}
    349 & 1\\
    174 & 0\\
    87 & 1\\
    43 & 1\\
    21 & 1\\
    10 & 0\\
    5 & 1\\
    2 & 0\\
    1 & 1\\
    0 & \\
  \end{tabular}}
  \end{minipage}\\[6pt]
\par\smallskip\noindent\hspace*{1em}\begin{tabular}{@{}p{0.42\linewidth}p{0.42\linewidth}@{}}
  \textbf{a)}~$478_{10} = 111011110_2$ & \textbf{b)}~$203_{10} = 11001011_2$\\[4pt]
  \textbf{c)}~$941_{10} = 1110101101_2$ & \textbf{d)}~$349_{10} = 101011101_2$
\end{tabular}\par\smallskip
  Si divide per 2 fino ad arrivare a 0 e si leggono i resti dal basso verso l'alto.}
  \criterioRR
  {2 pt per ogni numero: 1 pt divisioni per 2 corrette, 1 pt resti letti dal basso e risultato. Risultati: a) $10110011_2$; b) $110101011_2$; c) $100101110_2$; d) $1101100111_2$.}
  {2 pt per ogni numero: 1 pt divisioni per 2 corrette, 1 pt resti letti dal basso e risultato. Risultati: a) $111011110_2$; b) $11001011_2$; c) $1110101101_2$; d) $101011101_2$.}
\end{solution}
\end{esercizio}
